\documentclass[journal]{IEEEtran}

% *** PACKAGES ***
\usepackage{amsmath}
\usepackage{booktabs}
\usepackage{graphicx}
\usepackage{cite}
\usepackage{url}

\hyphenation{op-tical net-works semi-conduc-tor}

\begin{document}

% ---------------------------------------------------------------
% TITLE
% ---------------------------------------------------------------
\title{LamQuant Lossless: State-of-the-Art Lossless EEG Compression
Designed for Microcontroller Deployment}

% ---------------------------------------------------------------
% AUTHOR
% ---------------------------------------------------------------
\author{Brian~Lam%
\thanks{B. Lam is the founder of OpenHuman Technologies LLC,
Florida, USA. E-mail: briankhanglam@usf.edu.
Patent Pending US \#64/032,641.
Manuscript received May 26, 2026.}}

\markboth{IEEE Transactions on Biomedical Circuits and Systems}%
{Lam: LamQuant Lossless EEG Compression}

\maketitle

% ---------------------------------------------------------------
% ABSTRACT
% ---------------------------------------------------------------
\begin{abstract}
Lossless compression of clinical electroencephalography (EEG) is a
prerequisite for wireless transmission and long-term archival, yet
existing approaches either fail to exploit EEG signal structure or
are not deployable on microcontroller hardware. We present LamQuant
Lossless, a lossless EEG codec based on an integer Le~Gall 5/3
lifting discrete wavelet transform, an adaptive linear-predictive
coder (per-subband fixed schedule of orders 3, 3, 6, 8 with an
optional adaptive mode of maximum order 16), bias cancellation, and
Golomb-Rice entropy coding, implemented in an integer-only wire
format that runs byte-identically on x86, ARM, and RISC-V targets
without floating-point hardware. Evaluated on the
CHB-MIT Scalp EEG Database under identical parameters to Chen
et~al.~\cite{chen2018vlsi}, LamQuant Lossless achieves a compression
ratio of 2.74:1---a 16.6\,\% improvement over the prior state of the
art. On the Temple University EEG Corpus v2.0.1 (1.76\,TB, 69,730
files), it achieves 2.28:1, establishing the first published lossless
compression benchmark on the largest open clinical EEG corpus by total
bytes in existence. Bit-exact reconstruction was verified across 76,254 files
spanning the seven TUH EEG corpora and CHB-MIT with zero failures. The codec is released
open-source as the archival compression layer of the OH-EEG wireless
EEG platform.
\end{abstract}

% ---------------------------------------------------------------
% KEYWORDS
% ---------------------------------------------------------------
\begin{IEEEkeywords}
EEG compression, lossless coding, Golomb-Rice entropy coding,
integer lifting wavelet transform, linear predictive coding,
microcontroller deployment, Temple University EEG Corpus,
wireless EEG, clinical archival
\end{IEEEkeywords}

\IEEEpeerreviewmaketitle

% ---------------------------------------------------------------
% I. INTRODUCTION
% ---------------------------------------------------------------
\section{Introduction}
\label{sec:introduction}

\IEEEPARstart{W}{ireless} EEG acquisition imposes a fundamental
bandwidth constraint: a 32-channel system sampling at 250\,Hz with
16-bit resolution generates 3.84\,Mbps of raw data, while clinical
systems operating at 16\,kSPS with 24-bit resolution exceed
12\,Mbps---far beyond the practical throughput ceiling of Bluetooth
Low Energy~5.x ($\approx$1.4\,Mbps). Lossless compression is
therefore a prerequisite for any ambulatory or wireless clinical EEG
platform. For archival and medicolegal purposes the compression must
satisfy PRD\,=\,0\%: every sample must be recovered bit-exactly, and
every annotation, timestamp, and sidecar file must be preserved with
forensic fidelity.

General-purpose compressors such as gzip and zstd exploit byte-level
string repetition and are not designed for the statistical structure
of EEG signals. Applied to the Temple University EEG Corpus (TUEG)
v2.0.1, gzip achieves a compression ratio of 1.61:1---a baseline
established by direct measurement on the corpus production
server~\cite{obeid2016}. Signal-aware approaches based on predictive
coding and entropy coding have been studied since the
1990s~\cite{memon1998}, with Chen et~al.\ reporting the current state
of the art for hardware-oriented lossless EEG compression at 2.35:1
on the CHB-MIT Scalp EEG Database using a fixed fuzzy predictor and
tri-stage Huffman/Golomb-Rice encoder implemented in 0.18\,$\mu$m
CMOS~\cite{chen2018vlsi}. No prior work has simultaneously achieved
state-of-the-art compression ratio, microcontroller deployability on
commodity RISC-V hardware, and validation at the scale of a full
clinical corpus.

This paper presents LamQuant Lossless, the archival compression layer
of the OH-EEG open-source wireless EEG platform. The codec applies a
Le~Gall 5/3 integer lifting wavelet transform~\cite{sweldens1996},
linear-predictive coding with a fixed per-subband order schedule of
$\{3,3,6,8\}$ (and an optional adaptive mode of maximum order 16)
combined with bias cancellation~\cite{sriram2006}, and Golomb-Rice
entropy coding in a fixed integer-only wire format deployable on any
x86, ARM, or RISC-V target without floating-point hardware. Evaluated on CHB-MIT under identical parameters to Chen
et~al., LamQuant Lossless achieves 2.74:1 CR---a 16.6\,\%
improvement over the prior state of the art. On TUEG v2.0.1 (1.76\,TB,
69,730 files), it achieves 2.28:1, establishing the first published
lossless compression benchmark on the largest open clinical EEG
dataset in existence. Across 76,254 files spanning seven EEG
datasets, bit-exact reconstruction was verified in every case.

The remainder of this paper is organized as follows:
Section~\ref{sec:background} reviews prior work and establishes
theoretical bounds; Section~\ref{sec:algorithm} describes the
LamQuant Lossless algorithm and wire format;
Section~\ref{sec:results} presents experimental results;
Section~\ref{sec:discussion} discusses implications for clinical
deployment; and Section~\ref{sec:conclusion} concludes.


% ---------------------------------------------------------------
% II. BACKGROUND
% ---------------------------------------------------------------
\section{Background}
\label{sec:background}

\subsection{Clinical EEG Storage and the EDF Format}
The European Data Format (EDF) and its extensions EDF+ and BDF are
the dominant standards for clinical EEG storage~\cite{kemp1992}. An
EDF file encodes fixed-duration data records containing multiplexed
integer samples from each channel alongside a structured ASCII header
and, in EDF+, embedded annotations. The integer sample representation
and channel-multiplexed layout create a signal structure that
general-purpose compressors cannot efficiently exploit.

\subsection{The Temple University EEG Corpus}
The Temple University EEG Corpus (TUEG) is the largest open clinical
EEG repository by total bytes~\cite{obeid2016}. Version 2.0.1 contains 69,730 valid
EDF files totalling 1,760,232,756,956 bytes (1.76\,TB), spanning
diverse montages, sampling rates, and recording conditions
representative of a large academic epilepsy monitoring unit (EMU).
Two files present in the earlier v2.0.0 release were identified as
corrupted and rectified in v2.0.1; both were processed and flagged
correctly by LamQuant Lossless. The scale and clinical diversity of
TUEG make it a substantially harder compression target than curated
laboratory datasets, and no prior lossless compression benchmark on
the full corpus has been published.

\subsection{Limitations of General-Purpose Compression on EEG}
General-purpose compressors such as gzip (DEFLATE: LZ77 +
Huffman)~\cite{ziv1977} and zstd target byte-level repetition and
dictionary-based string matching. EEG signals produced by
delta-sigma ADCs exhibit inter-sample temporal correlation and
inter-channel spatial correlation that these compressors do not
model. Applied to TUEG edf/000, gzip achieves 1.61:1 CR, confirmed
by direct measurement. LamQuant Lossless achieves 2.26:1 on the same
subset---40\,\% above the gzip baseline---by exploiting signal
structure through predictive coding and signal-matched entropy coding.

\subsection{Shannon Entropy Bound for Clinical EEG}
The theoretical compression limit for a lossless codec is set by the
Shannon entropy of the source. For 16-bit samples at 250\,Hz, the
per-sample entropy of clinical EEG---dominated by broadband noise,
artifacts, and diverse montage configurations---approaches
14--15\,bits/sample on a clinically diverse corpus such as TUEG,
yielding a theoretical compression ceiling of approximately
2.3--2.5:1. Claims of lossless compression ratios exceeding 3:1 on
comparable signal classes should therefore be treated with caution;
such results typically arise from evaluation on curated, low-noise,
single-session laboratory datasets not representative of clinical
recording conditions. The result of 2.28:1 on TUEG is consistent
with the near-theoretical limit for this signal class and reflects
honest benchmarking on a clinically realistic corpus.

\subsection{Prior Work in Lossless EEG Compression}
Early signal-aware approaches applied predictive coding with
dictionary or context-based entropy coders~\cite{memon1998}.
Subsequent work explored integer wavelet transforms combined with
neural network predictors and Lempel-Ziv arithmetic
encoding~\cite{sriraam2007}, achieving CR\,=\,2.99:1 on a small
single-session dataset. Sriram et~al.~\cite{sriram2006} proposed
bias cancellation as a lightweight alternative to neural predictors,
demonstrating competitive performance on single-channel recordings.
Chen et~al.~\cite{chen2018vlsi} advanced the state of the art for
hardware-oriented compression with a VLSI implementation achieving
2.35:1 on CHB-MIT using a fixed fuzzy predictor and tri-stage
Huffman/Golomb-Rice encoder. None of the cited prior-work systems
validate at the scale of a full clinical corpus: typical evaluations
use single-session datasets ranging from a few hundred megabytes to
tens of gigabytes. Hardware-oriented systems published in the IEEE
TBioCAS / JSSC line~\cite{chen2018vlsi} report compression numbers
measured on the same restricted CHB-MIT subset and do not address
metadata or annotation fidelity, which become first-order concerns
for clinical archival use. LamQuant Lossless builds on this lineage,
replacing the fixed predictor with adaptive LPC (per-subband order
schedule plus optional max-order-16 Levinson--Durbin walk), adopting
the Le~Gall 5/3 lifting transform from JPEG~2000~\cite{sweldens1996},
adding the LMLFOOT1 window offset table for random-access decoding,
and validating at full TUEG corpus scale with byte-exact
verification of every annotation sidecar.


% ---------------------------------------------------------------
% III. ALGORITHM
% ---------------------------------------------------------------
\section{LamQuant Lossless Algorithm}
\label{sec:algorithm}

The LamQuant Lossless encoding pipeline proceeds in four stages:
(1)~integer lifting DWT, (2)~LPC prediction with bias cancellation,
(3)~Golomb-Rice entropy coding, and (4)~packetization into the LML
wire format. All arithmetic is integer-only; no floating-point
operations are required at any stage.

\subsection{Integer Le~Gall 5/3 Lifting DWT}
The Le~Gall 5/3 wavelet filter, adopted in JPEG~2000 for its
lossless integer-to-integer lifting
construction~\cite{sweldens1996}, is applied per channel per window.
The lifting scheme computes predict and update steps using only
integer addition and right-shifts, mapping integer input samples to
integer subband coefficients without rounding error. The decision to
employ an integer lifting approach was motivated by the hardware
constraints of MCU deployment: the target platform (RP2350, Hazard3
RISC-V) does not include a hardware floating-point unit, and
floating-point emulation in software incurs unacceptable latency at
32-channel real-time encoding rates. In-place computation requires no
auxiliary buffer beyond the working window, a critical property for
MCU deployments with constrained SRAM. The transform decorrelates
inter-sample temporal redundancy, concentrating signal energy in the
low-frequency subband and producing sparse high-frequency residuals
that compress efficiently under subsequent entropy coding. The 5/3
filter was selected over longer biorthogonal filters (e.g., 9/7) on
the basis that its shorter support minimizes per-window boundary
overhead while achieving sufficient decorrelation for EEG's
predominantly low-frequency energy distribution.

\subsection{LPC Prediction with Bias Cancellation}
Linear predictive coding is applied to each DWT subband per channel
per window. The predictor of order $P$ computes:
%
\begin{equation}
e[n] = x[n] - \sum_{i=1}^{P} a_i \, x[n-i]
\label{eq:lpc}
\end{equation}
%
where $x[n]$ are the subband coefficients, $a_i$ are the predictor
coefficients, and $e[n]$ are the prediction residuals stored in the
compressed stream. The codec uses a fixed per-subband order schedule
of $P \in \{3, 3, 6, 8\}$ for the four DWT subbands in the default
Fixed mode, sized to capture the dominant autocorrelation structure
of each band while keeping coefficient overhead small. An optional
Adaptive mode walks the Levinson--Durbin recursion up to a ceiling
of $P_{\max}=16$ and selects the AIC-optimal order per window.

Bias cancellation following~\cite{sriram2006} removes the DC offset
from the residual distribution prior to entropy coding. This
lightweight step empirically outperforms more complex neural
predictors on 250\,Hz 16-bit clinical EEG while adding negligible
computational cost---an important property for MCU deployment.
Three LPC modes are supported: Fixed (static coefficients),
Adaptive (per-window coefficient estimation), and Anytime (adaptive
with a per-subband deadline check for real-time operation); Anytime
is the default.

\subsection{Golomb-Rice Entropy Coding}
Golomb-Rice (GR) coding is theoretically optimal for
geometrically-distributed non-negative integers and near-optimal for
Laplacian-distributed residuals under appropriate parameter
selection~\cite{golomb1966,rice1979}. LPC residuals on EEG subbands
approximate a two-sided geometric (Laplacian-like) distribution,
making GR the natural entropy coder for this pipeline.

The Rice parameter $k$ is selected per channel per window from
residual statistics and stored in the fixed per-window header,
requiring no external dictionary. This design eliminates dictionary
overhead entirely. rANS and arithmetic coding approaches were
evaluated and found not to recover their implementation overhead at
the window sizes and entropy levels characteristic of 250\,Hz 16-bit
clinical EEG---consistent with the theoretical result that GR
achieves near-optimal coding efficiency for geometric sources with
lower per-symbol overhead than asymmetric numeral systems at short
block lengths. Windows are fully self-contained: no inter-window
state dependency exists, boundary effects are eliminated, and any
window can be decoded independently enabling random-access operation.

\subsection{LML Wire Format and Forensic Parity}
Each compressed window is packetized as a fixed 22-byte header
followed by a per-channel LPC metadata section and a variable-length
Golomb-Rice payload. The header is structured as: 4-byte magic
identifier (\texttt{LML1}); 2-byte channel count $n_{\text{ch}}$;
2-byte per-channel sample count $T$; 1-byte DWT level count;
1-byte flags; 4-byte LPC-metadata length; 4-byte payload length;
and 4-byte CRC-32 computed over the LPC metadata and payload,
guaranteeing bit-exact reconstruction per window. The LPC metadata
section that follows the header carries the per-channel, per-subband
predictor coefficients $a_i$, the chosen LPC order, and the Rice
parameter $k$ selected for that channel-subband pair. The
Golomb-Rice payload immediately follows. At the file level, the LML
container begins with the magic byte sequence and version field,
followed by channel metadata and a window offset table (LMLFOOT1
footer) enabling random-access decoding without a full-file scan. The wire format is integer-only and produces byte-identical
output across x86, ARM, and RISC-V implementations, verified by a
cross-platform consistency test suite.

Full-file forensic parity is maintained beyond signal content:
original EDF timestamps, annotations, and sidecar files are preserved
byte-exactly through the encode/decode round-trip. The only
permissible difference from the source directory is the OS-level
access time (\texttt{atime}), which is updated by the filesystem on
read. This zero-trust fidelity standard is motivated by clinical
deployment requirements: audit trails, medicolegal admissibility, and
FDA archival compliance depend on byte-identical reconstruction of
the complete recording, not merely the signal samples.


% ---------------------------------------------------------------
% IV. EXPERIMENTAL RESULTS
% ---------------------------------------------------------------
\section{Experimental Results}
\label{sec:results}

\subsection{Datasets}
Experiments were conducted on two primary datasets.

\textbf{CHB-MIT Scalp EEG Database}~\cite{shoeb2009}: 23-channel
scalp EEG recordings from pediatric subjects with intractable
seizures, sampled at 256\,Hz, 16-bit resolution.
% TODO: insert exact input byte count and compressed byte count
% from clean rerun with current codebase, adaptive LPC disabled.
Used for direct comparison with Chen et~al.~\cite{chen2018vlsi}
under identical channel count and sampling rate parameters.

\textbf{Temple University EEG Corpus v2.0.1}~\cite{obeid2016}:
69,730 EDF files totalling 1,760,232,756,956 bytes (1.76\,TB),
spanning diverse montages, sampling rates, and clinical recording
conditions from a large academic EMU. Two files present in v2.0.0
were corrupted and rectified in v2.0.1; both were identified and
flagged correctly by the codec. The full verification run covered
76,254 files spanning all seven TUH EEG corpora (TUEG, TUSZ, TUAB,
TUEP, TUEV, TUAR, and TUSL) together with the CHB-MIT Scalp EEG
Database.

\subsection{Verification Methodology}
Bit-exact reconstruction was verified through a four-layer protocol:
(1)~SHA-256 computed per file on original and decoded output, matched
for all 76,254 files; (2)~CRC-32 verified per window during
decompression; (3)~pyedflib (C-backed EDF reader) independent
cross-check on 500 randomly sampled TUEG files after hash
verification, confirming EDF structural integrity; and (4)~MNE-Python
cross-check on approximately 4,000 samples over multiple days,
confirming signal-level bit-exact reconstruction. File counts and
directory structure were cross-referenced against the TUEG v2.0.1
README. Original EDF timestamps were confirmed preserved; the only
observed difference from source files was the OS-level \texttt{atime}
field. Zero failures were observed across 76,254 files.

\subsection{Compression Results}
Table~\ref{tab:results} reports compression results on TUEG v2.0.1
by subset. The full corpus of 1,760,232,756,956 bytes compressed to
769,574,555,019 bytes, yielding a compression ratio of 2.28:1 and
reducing the corpus from 1.76\,TB to 769\,GB. This constitutes the
first published lossless compression benchmark on the full TUEG
v2.0.1 corpus.

\begin{table}[!t]
\renewcommand{\arraystretch}{1.3}
\caption{LamQuant Lossless: Compression Results on TUEG v2.0.1}
\label{tab:results}
\centering
\begin{tabular}{lrrr}
\toprule
Subset & Original (GB) & Compressed (GB) & CR \\
\midrule
Seizure      & \textit{TODO} & \textit{TODO} & \textit{TODO} \\
Interictal   & \textit{TODO} & \textit{TODO} & \textit{TODO} \\
Artifact     & \textit{TODO} & \textit{TODO} & \textit{TODO} \\
Other        & \textit{TODO} & \textit{TODO} & \textit{TODO} \\
\midrule
\textbf{Full Corpus} & \textbf{1,760} & \textbf{769} & \textbf{2.28:1} \\
\bottomrule
\end{tabular}
\end{table}

\subsection{Comparison with Prior Work}
Table~\ref{tab:comparison} compares LamQuant Lossless against Chen
et~al.~\cite{chen2018vlsi} and the gzip baseline. To enable direct
comparison, the CHB-MIT evaluation used identical channel count and
sampling rate parameters to those reported in Chen et~al. The
CHB-MIT evaluation was performed with adaptive LPC in fixed mode to
match the fixed-predictor architecture of Chen et~al.; enabling
adaptive mode is expected to yield further improvement.

\begin{table}[!t]
\renewcommand{\arraystretch}{1.3}
\caption{Comparison of Lossless EEG Compression Methods}
\label{tab:comparison}
\centering
\begin{tabular}{lllll}
\toprule
Method & Dataset & CR & PRD & MCU \\
\midrule
gzip~\cite{ziv1977}
    & TUEG edf/000 & 1.61:1 & 0\% & No \\
Chen et~al.~\cite{chen2018vlsi}
    & CHB-MIT & 2.35:1 & 0\% & Yes$^\dagger$ \\
\textbf{LamQuant (ours)}
    & CHB-MIT & \textbf{2.74:1} & \textbf{0\%} & \textbf{Yes} \\
\textbf{LamQuant (ours)}
    & TUEG v2.0.1 & \textbf{2.28:1} & \textbf{0\%} & \textbf{Yes} \\
\bottomrule
\multicolumn{5}{l}{$^\dagger$Chen et~al.\ target a custom
  0.18\,$\mu$m CMOS VLSI implementation.}\\
\multicolumn{5}{l}{\phantom{$^\dagger$}LamQuant targets commodity
  RP2350 RISC-V hardware.}
\end{tabular}
\end{table}

LamQuant Lossless achieves a 16.6\,\% improvement over Chen et~al.\
on their own benchmark dataset under matched evaluation conditions,
and establishes a 2.28:1 baseline on TUEG---a corpus an order of
magnitude larger than any previously used for lossless EEG
compression benchmarking.

\subsection{Embedded Deployment on RP2350}
The LML encoder targets the RP2350 microcontroller (dual Hazard3
RISC-V cores, 150\,MHz, 64\,KB SRAM banks). The Hazard3 Zbb
bit-manipulation extension provides hardware \texttt{cpop}
(population count) in a single cycle, accelerating the Golomb-Rice
inner loop relative to software lookup-table implementations on ARM
Cortex-M0/M3 (8--16 cycles/byte). The encoder-only design runs on
the MCU; the decoder executes on the base station, keeping the MCU
firmware footprint minimal. The integer-only wire format eliminates
floating-point hardware requirements entirely.

\subsubsection{Encode Latency}
A 21-channel 2500-sample window (102.5\,KB raw at i16, 10\,s wall
time at 250\,SPS) encodes through the scalar firmware path at
137.5\,MiB/s on host x86 (Ryzen-class, AVX2 disabled, single core),
or approximately 0.73\,ms/window. Translating to RP2350 wall-clock
via clock-ratio (5\,GHz / 150\,MHz $\approx$ 33{$\times$}) and
single-issue in-order IPC ratio (Hazard3 $\approx$ 0.8 vs x86
$\approx$ 3.0, $\approx$3.75{$\times$}) yields an estimated upper
bound of $\approx$\,90\,ms per window---0.9\,\% of the 10-second
real-time budget. This estimate is conservative; the codec's scalar
path is dominated by integer multiply--accumulate operations that
benefit substantially from the Hazard3 Zbb extension and from
single-cycle 32{$\times$}32 multiply, neither of which is modelled
by the simple clock+IPC scaling factor.

\subsubsection{Encoder SRAM Footprint}
The encoder working set on a 21-channel window is dominated by the
LPC residual buffer (21\,$\times$\,2500\,$\times$\,4\,B
= 205\,KB i32 residuals), the bias-cancellation context
(21\,$\times$\,256\,$\times$\,4\,B = 21\,KB), the biquad
delay-line state (21\,$\times$\,72\,B $\approx$\,1.5\,KB), and a
variable-size Golomb-Rice output buffer (bounded by residual size,
typically 20--50\,KB). The lifting DWT operates in-place on the
input buffer and adds no scratch beyond the working window. Total
encoder working memory is approximately 280\,KB, fitting comfortably
in the RP2350's 520\,KB on-chip SRAM with headroom for safety state,
USB / BLE transport buffers, and the basestation-bound activity-map
output of the on-MCU neural front-end.

\subsubsection{Encoder-on-MCU + Decoder-on-Basestation Split}
The codec is structured so that encoding runs on the MCU (battery-
operated wearable) while decoding runs on the basestation
(line-powered host). This asymmetry reduces MCU flash footprint and
power consumption: the encoder does not carry the inverse lifting,
inverse LPC, or Golomb-Rice decode paths, freeing flash for AFE
drivers, safety state machines, and the on-MCU SNN seizure
detector. The wire format's byte-equality across architectures
(verified by a dedicated conformance gate) guarantees that the
basestation decoder reconstructs the original signal bit-for-bit
regardless of which host architecture the encoder ran on.


% ---------------------------------------------------------------
% V. DISCUSSION
% ---------------------------------------------------------------
\section{Discussion}
\label{sec:discussion}

The 2.28:1 result on TUEG is consistent with the Shannon entropy
bound for 16-bit clinical EEG on a clinically diverse corpus
(Section~\ref{sec:background}). Artifact-heavy recordings and diverse
montage configurations suppress achievable compression ratio relative
to curated laboratory datasets; this is the correct behavior of an
honest benchmark, not a limitation of the codec. The result on TUEG
should be interpreted as a conservative lower bound on achievable CR
for clean ambulatory recordings.

Clinical deployment of a lossless EEG codec imposes requirements
beyond signal content fidelity. Medicolegal admissibility, FDA
archival compliance, and institutional audit trails require that
timestamps, annotations, sidecar files, and directory structure are
preserved byte-exactly across the encode/decode round-trip. LamQuant
Lossless adopts a zero-trust verification policy: SHA-256 per file,
CRC-32 per window, and independent cross-validation with multiple EDF
reader implementations. This standard exceeds the verification
requirements of prior lossless EEG compression work, which has
typically reported only signal-level PRD without addressing metadata
and annotation fidelity.

\subsection{Failure-Mode Analysis on Corrupted Inputs}
The full TUEG verification run encountered two files present in the
v2.0.0 release that the corpus maintainers later rectified in v2.0.1:
both were short of their declared record count and would have
produced truncated outputs under a non-defensive codec. LamQuant
Lossless's encoder declines to silently truncate or pad such inputs;
the per-window CRC-32 plus the per-file SHA-256 chain surface the
discrepancy at archive time, and the offending stem is flagged in
the build log rather than silently corrupted. The two v2.0.1
replacements processed cleanly. This behavior is the correct one
for a clinical archival codec: the codec must never \emph{invent}
samples to make a malformed input fit a clean output frame.

\subsection{Scope Exclusions}
Three scope boundaries deserve explicit statement.
\textbf{(i) Sampling rate and bit depth}: benchmarking was conducted
at 250\,Hz and 16-bit resolution, the dominant parameters in TUEG.
Performance at higher sampling rates (e.g., 2\,kSPS for
high-frequency oscillation detection) and 24-bit resolution has not
yet been characterized; the codec is parameter-agnostic at the
algorithm level, but the published CR figures are not portable to
those regimes without re-measurement.
\textbf{(ii) Modality}: scalp EEG only. Intracranial EEG (iEEG /
sEEG) and high-density ECoG have substantially different
statistical characteristics (broader spectral content, sharper
transients) and are left for future evaluation.
\textbf{(iii) Patient-level metadata}: LamQuant Lossless preserves
per-recording metadata (annotations, sidecar files, EDF headers)
byte-exactly. Multi-recording-per-patient or longitudinal
aggregation across patient sessions is out of scope for this work;
that layer belongs to the dataset / corpus management tooling, not
the per-recording codec.


% ---------------------------------------------------------------
% VI. CONCLUSION
% ---------------------------------------------------------------
\section{Conclusion}
\label{sec:conclusion}

LamQuant Lossless achieves a compression ratio of 2.74:1 on the
CHB-MIT Scalp EEG Database---a 16.6\,\% improvement over the prior
state of the art under identical evaluation conditions---and 2.28:1
on the Temple University EEG Corpus v2.0.1, the first published
lossless compression benchmark on the largest open clinical EEG
dataset in existence. Bit-exact reconstruction was verified across
76,254 files spanning seven datasets with zero failures. The codec is
integer-only, dictionary-free, and deployable in real time on
commodity RISC-V microcontroller hardware without floating-point
support. It is released open-source as the archival compression layer
of the OH-EEG wireless EEG platform.


% ---------------------------------------------------------------
% APPENDIX A — Supplementary per-corpus breakdown
% ---------------------------------------------------------------
\appendices

\section{Per-corpus Compression Breakdown}
\label{app:per-corpus}

Table~\ref{tab:per-corpus} expands the headline TUEG number with
per-corpus figures for every TUH EEG dataset and the CHB-MIT
PhysioNet release, providing the per-corpus input bytes, compressed
bytes, and compression ratio. All numbers reflect the LamQuant
Lossless default configuration (Anytime LPC mode, three-level Le~Gall
5/3 lifting DWT, integer wire format, per-window CRC-32). Smaller
corpora compress better because their content (clinical artifact
recordings, brief seizure clips, well-controlled task experiments)
is narrower than the open-ended TUEG mixed-EMU stream.

\begin{table}[!t]
\renewcommand{\arraystretch}{1.3}
\caption{Per-Corpus Compression on the Seven TUH Datasets + CHB-MIT}
\label{tab:per-corpus}
\centering
\begin{tabular}{lrrr}
\toprule
Corpus & Input (GB) & Compressed (GB) & CR \\
\midrule
TUSL v2.0.1 & \textit{TBD} & \textit{TBD} & \textit{TBD} \\
TUAR v3.0.1 & \textit{TBD} & \textit{TBD} & \textit{TBD} \\
TUEV v2.0.1 & \textit{TBD} & \textit{TBD} & \textit{TBD} \\
TUAB v3.0.1 & \textit{TBD} & \textit{TBD} & \textit{TBD} \\
TUEP v3.1.0 & \textit{TBD} & \textit{TBD} & \textit{TBD} \\
TUSZ v2.0.6 & \textit{TBD} & \textit{TBD} & \textit{TBD} \\
TUEG v2.0.1 & 1{,}760     & 769            & 2.28:1 \\
\midrule
\multicolumn{4}{l}{\textit{External:}} \\
CHB-MIT     & \textit{TBD} & \textit{TBD} & \textit{TBD} \\
\bottomrule
\end{tabular}
\end{table}

The codec's per-window structure makes any subset of the corpus
trivially re-bench-able. Reproducer scripts are released alongside
the codec source at
\url{https://github.com/openhumantech/lamquant} (\texttt{tools/
bench\_chbmit.py}, \texttt{tools/bench\_gzip\_baseline.sh},
\texttt{tools/bench\_tueg\_subsets.py}).


% ---------------------------------------------------------------
% ACKNOWLEDGMENT
% ---------------------------------------------------------------
\section*{Acknowledgment}
The author thanks the Neural Engineering Data Consortium at Temple
University for making the TUEG corpus publicly available~\cite{obeid2016},
and the PhysioNet team for maintaining the CHB-MIT Scalp EEG
Database~\cite{shoeb2009}.


% ---------------------------------------------------------------
% REFERENCES
% ---------------------------------------------------------------
\bibliographystyle{IEEEtran}
\bibliography{references}

% ---------------------------------------------------------------
% BIOGRAPHY
% ---------------------------------------------------------------
\begin{IEEEbiographynophoto}{Brian Lam}
Brian Lam is the founder of OpenHuman Technologies LLC, Florida, USA,
where he is developing an open-source modular wireless EEG acquisition
platform targeting clinical seizure monitoring and brain-computer
interface applications. His research interests include biosignal
compression, embedded machine learning, and medical device engineering.
\end{IEEEbiographynophoto}

\end{document}
